\documentclass[11pt,a4paper]{ctexart}

% =========================================================
% Packages
% =========================================================
\usepackage[a4paper,margin=2.4cm]{geometry}
\usepackage{amsmath,amssymb,amsthm}
\usepackage{bm}
\usepackage{booktabs}
\usepackage{array}
\usepackage{multirow}
\usepackage{xcolor}
\usepackage{hyperref}
\usepackage{tikz}
\usepackage{enumitem}
\usepackage{microtype}

\usetikzlibrary{
    positioning,
    arrows.meta,
    shapes.geometric,
    calc
}

\hypersetup{
    colorlinks=true,
    linkcolor=blue!60!black,
    citecolor=blue!60!black,
    urlcolor=blue!60!black
}

% =========================================================
% Custom commands
% =========================================================
\newcommand{\DeepCAVE}{DeepCAVE}
\newcommand{\DeepCAVEX}{DeepCAVE-X}
\newcommand{\HPO}{HPO}

\newcommand{\todo}[1]{%
    \textcolor{red}{\textbf{[REF/TO-DO: #1]}}%
}

% =========================================================
% Document
% =========================================================
\begin{document}

\title{\textbf{DeepCAVE-X: From HPO Visualization to Decision-Centric Optimization Analysis}}
\author{}
\date{}

\maketitle

\begin{abstract}
超参数优化（HPO）已经从一次性试错发展为基于历史观测更新搜索策略、资源分配和停止规则的序贯决策过程，但现有可视分析工具多将 HPO run 呈现为若干相互独立的终态视图，难以回答“当前结果是否足够、为什么停滞、应如何干预”等决策问题。本文提出 \DeepCAVEX{}，一个以 decision loop 为中心的 HPO 可视分析扩展框架。该方法将 HPO 分析组织为 Monitor、Attribute、Diagnose、Intervene 和 Re-optimize 五个阶段，并基于 DeepCAVE 的插件体系实现 \textbf{Weighted Sobol}、\textbf{Sampling Density}、\textbf{Footprint Replay}、\textbf{Bottleneck Diagnosis}、\textbf{Hypervolume Convergence}、\textbf{Fidelity Replay} 与 \textbf{Web Progress Feedback}。其中，Weighted Sobol 将全局敏感性分析中的均匀测度替换为优化器经验采样测度；Sampling Density 揭示重要维度上的边际与联合覆盖；Footprint Replay 为配置投影补充首次访问和 incumbent 迁移时间轴；Bottleneck Diagnosis 结合改进率、分布检验和新颖率区分不同停滞模式；Hypervolume Convergence 将多目标前沿转化为 anytime 收敛证据；Fidelity Replay 重放 multi-fidelity 调度以量化“省多少算力、亏多少最优性”。由此，DeepCAVE-X 把可视化从优化结束后的结果解释，扩展为支持停止、继续或调整搜索空间与 fidelity 策略的决策接口。
\end{abstract}

% =========================================================
\section{设计思路}

超参数优化（Hyperparameter Optimization, HPO）的目标是在给定配置空间和评估预算内寻找性能最优的超参数配置。随着 Bayesian Optimization、Successive Halving、Hyperband 及其他自动化搜索方法的发展，HPO 已经从一次性试错演变为基于历史观测的序贯决策过程：优化器根据已有配置--性能观测更新 surrogate、acquisition function 或资源分配规则，并决定下一个配置与 fidelity \cite{jones1998efficient,snoek2012practical,li2018hyperband,falkner2018bohb,bischl2021hpo}。这种序贯性也使停止规则和预算分配成为 HPO 方法设计的核心问题，而不只是搜索结束后的附带分析 \cite{li2018hyperband,falkner2018bohb,holzleitner2021hpo}。因此，一次 HPO run 并不只是一个待排序的 trial 集合，而是一个不断演化的决策历史。用户在优化过程中通常需要反复回答若干相互关联的问题：当前结果是否已经足够？优化器是在有效收敛还是陷入停滞？哪些超参数真正影响目标？搜索空间是否被过度收缩或无效探索？下一轮应该继续、停止，还是调整搜索空间、采样器或 fidelity 策略？这些问题的答案共同构成了 HPO 的 \emph{decision loop} \cite{bergstra2012random,hutter2014importance}。

可视分析研究将这一过程概括为“分析、推理、决策与行动”的闭环：系统把数据转化为交互式证据，用户依据证据形成判断并触发后续分析 \cite{keim2008visual}。HyperTuner 等系统已经表明，交互式可视化能够帮助用户理解超参数、性能和收敛过程之间的关系 \cite{hyptertuner2018}。DeepCAVE 进一步提供了较为完整的 HPO 分析平台，支持 cost over time、configuration footprint、hyperparameter importance、partial dependence、Pareto front 和 budget correlation 等视图，并具有插件化架构和运行中分析能力 \cite{segel2025deepcave}。然而，从决策闭环的角度看，这些视图仍然主要是\emph{描述性}的：它们分别展示某一段时间内的性能轨迹、终态搜索覆盖或单参数效应，但没有显式回答“该证据应触发什么干预”。这对应可视分析中从 evidence generation 到 decision support 的关键过渡：仅有 insight 并不等于 action，用户仍需在多个插件之间自行拼接结论。更重要的是，现有视图在若干关键环节存在结构性缺口：
\begin{itemize}
    \item \textbf{性能评估不完整。} Cost over Time 能展示单目标 incumbent 曲线，Pareto Front 能展示静态非支配解，但二者难以直接回答“优化是否已经收敛、是否值得继续”。HPO 评估中已经广泛使用 anytime performance 来刻画有限预算下的质量演化，用 regret 衡量与已知最优或事后最优的差距 \cite{erickson2020autogluon,klein2017fast}；多目标优化也以 hypervolume 作为标准集合质量指标 \cite{zitzler1999spea,while2006faster}。但多目标场景缺少 hypervolume 随时间演化的统一收敛视图；单目标场景也只能看到曲线平台，而无法判断平台是正常收敛、局部停滞还是无效探索。
    \item \textbf{重要性与实际搜索分布脱节。} fANOVA 基于随机森林 surrogate 将性能方差分解为单参数和低阶交互贡献，LPI 则度量 incumbent 附近扰动的局部效应 \cite{hutter2014importance,biedenkapp2018cave}。Sobol 方差分解进一步提供了一阶效应、总效应和交互贡献的标准定义 \cite{sobol2001global,saltelli2010variance}，且敏感性分析的结果依赖于输入概率测度 \cite{kucherenko2012dependent}。HPO 优化器产生的样本通常高度非均匀，因此只依赖均匀或全局代理测度的重要性可能夸大未探索区域的作用，难以反映当前搜索分布下的实际可干预维度。
    \item \textbf{搜索行为缺少过程维度。} Bayesian optimization 的 acquisition 决策具有路径依赖性，探索--利用权衡会直接改变后续样本分布 \cite{snoek2012practical,frazier2018tutorial,shahriari2016taking}。Configuration Footprint 将配置空间压缩为终态 2D 覆盖图，能看出“搜索到了哪里”，但无法看出“何时开始收缩、哪些点先被访问、incumbent 如何迁移”。MDS 的作用是保持高维配置距离的可视近似 \cite{kruskal1964mds}，但冻结的终态投影缺少过程时间轴；采样密度也没有被显式展示，用户难以判断某个重要维度是被充分探索、过度集中还是贴边受限。
    \item \textbf{资源与 fidelity 决策证据不足。} Successive Halving、Hyperband、BOHB 和 ASHA 表明，fidelity 与预算分配本身就是 HPO 的序贯决策对象 \cite{jamieson2016nonstochastic,li2018hyperband,falkner2018bohb}；Budget Correlation 可以检查低预算与高预算结果的秩相关，但不能反事实地比较不同 multi-fidelity 调度可能带来的节省和性能损失。同时，静态插件长时间计算时缺少阶段化进度反馈，降低了用户在交互闭环中的可控性。
    \item \textbf{视图与干预动作之间缺少显式映射。} 用户可以看到重要性、覆盖和曲线，但系统没有把“重要性高且采样集中”“采样分散但无改进”“低 fidelity 不可靠”等证据组合成可操作的下一轮建议。
\end{itemize}

基于上述观察，本文将 HPO 分析组织为一个以决策为中心的迭代闭环：
\begin{equation}
\text{Monitor}
\rightarrow
\text{Attribute}
\rightarrow
\text{Diagnose}
\rightarrow
\text{Intervene}
\rightarrow
\text{Re-optimize}.
\label{eq:loop}
\end{equation}
其中，\textbf{Monitor} 评估当前性能是否已经足够；\textbf{Attribute} 识别哪些超参数值得进一步分析；\textbf{Diagnose} 结合搜索行为与资源使用解释性能瓶颈；\textbf{Intervene} 将诊断结果转化为对搜索空间、采样策略或 fidelity 的调整；\textbf{Re-optimize} 通过新的 run 或扩展后的 run 验证干预效果。该框架并不试图自动替代优化器或用户决策，而是把原本分散的视图组织成能够支持迭代干预的分析路径；这与人本交互式 HPO 研究的出发点一致，即让用户在多目标或中间反馈场景中参与优化方向的塑造 \cite{giovanelli2023interactive}。

\DeepCAVEX{} 的插件设计正是围绕这一闭环展开。表~\ref{tab:plugin-loop} 将现有插件和新增插件映射到五个阶段：表中新增插件以粗体标出，用于补足 DeepCAVE 原有视图从描述到干预的缺口。

\begin{itemize}
    \item \textbf{Hypervolume over Time} 补全性能监控环节。Pareto dominance 只给出支配关系，Hypervolume 则以参考点为界度量非支配集合所覆盖的目标空间 \cite{zitzler1999spea}，已被广泛用于多目标算法收敛性与选择压力评估 \cite{while2006faster}。将其绘制为时间或 trial 序列，可以把静态前沿转化为 anytime convergence 证据，从而支持“继续优化还是停止”的判断。
    \item \textbf{Bottleneck Detection} 将性能平台转化为可诊断事件。BO 的探索--利用理论表明，无改进既可能来自过度开采，也可能来自无效探索 \cite{srinivas2010gaussian,shahriari2016taking}。该插件因此结合 incumbent 相对改进和搜索探索信号，区分局部最优型停滞与无效探索型停滞，避免用户把所有平台都误读为“应当停止”，也避免在没有证据的情况下盲目追加预算。
    \item \textbf{Sobol / Weighted Sobol} 补强归因环节。Sobol 方差分解将输出不确定性分配到输入子集，并给出标准的一阶指数与总效应指数 \cite{sobol2001global,saltelli2010variance}；测度改变会改变敏感性解释，这在依赖非均匀输入的敏感性研究中已被明确讨论 \cite{kucherenko2012dependent}。Weighted Sobol 进一步以优化器实际经验采样分布替代均匀测度，回答“在当前搜索已经覆盖的区域中，哪些参数对性能和下一步干预更重要”。因此，它不是简单替换 fANOVA，而是把重要性分析与搜索行为联系起来。
    \item \textbf{Sampling Density} 展示重要维度上的实际采样分布。一维边际密度和二维联合密度对应敏感性分析中的 marginal 与 joint sampling behavior \cite{saltelli2008global}，可以揭示过度集中、贴边、均匀覆盖不足或类别偏好等模式，从而帮助用户决定应扩大、收缩还是保持某个超参数的搜索范围。
    \item \textbf{Footprint Replay} 为终态 footprint 增加时间轴。MDS 是保持样本间不相似性的经典低维投影方法 \cite{kruskal1964mds}。通过冻结 MDS 坐标并按 trial 回放首次访问点和 incumbent 前缀，用户可以观察搜索从探索到收缩的演化过程，判断优化器是否过早收敛、后期是否仍在大范围无效扩散。
    \item \textbf{Fidelity Replay} 把 fidelity 分析从静态相关性扩展到反事实调度比较。Successive Halving 与 Hyperband 的核心是通过保守低 fidelity 筛选和激进高 fidelity 提升，把有限预算分配给更有希望的配置 \cite{jamieson2016nonstochastic,li2018hyperband}；ASHA 进一步给出异步推广形式 \cite{li2020asha}。该插件回放这些调度，估计节省时间、incumbent 差距和被剪枝 trial 数，为“是否更换 multi-fidelity 策略”提供量化证据。
    \item \textbf{Web Progress Feedback} 支撑人在回路中的交互体验。可视分析系统需要在用户、算法与界面之间维持可感知的交互状态 \cite{keim2008visual,endert2012human}。静态插件计算可能持续数十秒到数分钟，阶段化进度反馈使用户理解系统当前状态，减少长时间等待造成的中断和误操作，从而维持分析--决策闭环的可用性。
\end{itemize}

综上，\DeepCAVEX{} 的贡献不是若干孤立插件的叠加，而是把 HPO run 重新组织为可诊断、可干预、可验证的决策过程。用户可以沿“性能监控 $\rightarrow$ 参数归因 $\rightarrow$ 搜索行为诊断 $\rightarrow$ 空间或 fidelity 干预 $\rightarrow$ 新一轮验证”的路径使用这些插件，并在每个阶段优先选择表~\ref{tab:plugin-loop} 中对应的证据视图。这样，可视化不再只是优化结束后的解释工具，而成为 HPO 迭代过程中的决策接口，也把 visual analytics 的分析闭环具体化为 HPO 的诊断与干预闭环 \cite{keim2008visual,sach2016progress}。
\begin{figure}[t]
    \centering
    \begin{tikzpicture}[
        node distance=0.7cm and 0.45cm,
        block/.style={
            draw,
            rounded corners,
            align=center,
            minimum width=2.25cm,
            minimum height=0.9cm,
            font=\small
        },
        arrow/.style={->, thick}
    ]
        \node[block] (perf) {Monitor\\
        当前结果是否足够？};
        \node[block, right=of perf] (attr) {Attribute\\
        哪些参数重要？};
        \node[block, right=of attr] (diag) {Diagnose\\
        搜索为何停滞？};
        \node[block, below=of diag] (inter) {Intervene\\
        下一步如何调整？};
        \node[block, left=of inter] (reo) {Re-optimize\\
        验证干预效果？};

        \draw[arrow] (perf) -- (attr);
        \draw[arrow] (attr) -- (diag);
        \draw[arrow] (diag) -- (inter);
        \draw[arrow] (inter) -- (reo);
        \draw[arrow] (reo) -| ([xshift=-0.25cm]perf.south);

        \node[below=0.25cm of reo, align=center, font=\small]
        {重新运行 HPO，形成下一轮决策闭环};
    \end{tikzpicture}
    \caption{Decision-centric workflow of DeepCAVE-X.}
    \label{fig:workflow}
\end{figure}

\begin{table}[t]
    \centering
    \caption{\DeepCAVE{} 与 \DeepCAVEX{} 插件在 decision loop 中的映射。新增插件以粗体表示。}
    \label{tab:plugin-loop}
    \small
    \begin{tabular}{@{}p{2.2cm}p{5.3cm}p{5.8cm}@{}}
        \toprule
        \textbf{Loop 阶段} & \textbf{DeepCAVE 现有插件} & \textbf{\DeepCAVEX{} 新增插件} \\
        \midrule
        Monitor & Cost over Time; Pareto Front & \textbf{Hypervolume over Time} \\
        Attribute & Hyperparameter Importance & \textbf{Sobol / Weighted Sobol} \\
        Diagnose & Configuration Footprint & \textbf{Bottleneck Detection}; \textbf{Sampling Density}; \textbf{Footprint Replay} \\
        Intervene & Budget Correlation & \textbf{Fidelity Replay}; \textbf{Sampling Density} \\
        Re-optimize & 复用 Monitor 视图 & 复用 Monitor 与新增收敛证据 \\
        \midrule
        全局支撑 & Plugin framework & \textbf{Web Progress Feedback} \\
        \bottomrule
    \end{tabular}
\end{table}


% =========================================================
\section{实现思路}
% =========================================================

\label{sec:implementation}

\DeepCAVEX{} 不修改 \DeepCAVE{} 的主框架，而是遵循其官方插件开发指导，在既有插件体系中实现新功能。具体而言，系统继续以 Run 作为统一数据模型，通过 converters 复用 SMAC、BOHB、Optuna、Ray Tune 等优化器日志的解析能力；新插件继承 StaticPlugin 或 DynamicPlugin 的接口与生命周期，并继续通过 Redis Queue 分发耗时计算和缓存结果。前端则沿用 Dash 提供的 Web GUI 与 API 双入口。整体实现架构如图~\ref{fig:architecture} 所示。

\begin{figure}[t]
    \centering
    \includegraphics[width=\textwidth]{interface/deepcave-structure.png}
    \caption{\DeepCAVE{} 的实现架构：Dash 提供 Web 与 API 双入口，converters 将不同优化器日志统一为 Run 对象，plugins 分为 Static（Redis Queue 异步计算）与 Dynamic（实时刷新）两类。}
    \label{fig:architecture}
\end{figure}

这种基于插件体系的扩展方式具有三方面好处：首先，它保证新功能能够与现有插件共存，不破坏已有数据读取、缓存和交互机制；其次，算法层 evaluator 与展示层 plugin 分离，使每个改进点可以独立开发、测试和复用；最后，用户可以在同一界面中将新旧插件组合成完整的 HPO 决策流程，而不需要在多个分析工具之间切换。


% =========================================================
\section{插件实现}
% =========================================================

本节说明 DeepCAVE-X 中三个核心插件的实现方式。它们分别服务于 Attribute、Diagnose 与 Intervene 阶段，并延续第~\ref{sec:implementation} 节中的 evaluator--plugin 分层结构。

\subsection{Weighted Sobol Importance}
\label{sec:weighted-sobol}

DeepCAVE 原有 Hyperparameter Importance 主要基于均匀输入测度下的全局方差分解，但 HPO 优化器产生的样本通常高度非均匀。为刻画``当前搜索已经覆盖的区域中，哪些超参数真正影响性能''，Weighted Sobol 不直接改写重要性数值，而是把 Sobol 方差分解中的测度由均匀分布替换为优化器的经验采样测度。该思路依赖 Kucherenko 等人对相关/非均匀输入下全局敏感性指标的推广：敏感性指数不再假设输入独立同均匀分布，而是相对于指定概率测度 $q(\mathbf{x})$ 计算~\cite{kucherenko2012dependent}。

设 $\mathbf{x}=(x_1,\ldots,x_d)$，经验测度 $q(\mathbf{x})$ 由各维边际分布的乘积近似。对连续超参数，经验边际采用 bootstrap 加核带宽 $1.06\,\sigma n^{-1/5}$ 的扰动；类别超参数按观测频率重采样。用一个随机森林回归器拟合成功 trial 的配置--性能样本，记预测响应为 $f(\mathbf{x})$。给定测度 $q$ 时，先估计总方差
\begin{equation}
D_q
=
\operatorname{Var}_{q}\!\left[f(\mathbf{x})\right]
=
\int f(\mathbf{x})^2 q(\mathbf{x})\,d\mathbf{x}
-
\left(\int f(\mathbf{x}) q(\mathbf{x})\,d\mathbf{x}\right)^2 .
\label{eq:weighted-variance}
\end{equation}
按 Kucherenko 等给出的相关/非均匀输入敏感性定义，第 $i$ 个超参数的一阶贡献可由替换第 $i$ 列得到的模型响应估计~\cite{kucherenko2012dependent}：
\begin{equation}
S_i^{q}
=
\frac{1}{D_q}
\int f(\mathbf{x}) f(\mathbf{x}_{\sim i}, z_i)
\, q(\mathbf{x}) q(z_i)\,d z_i\,d\mathbf{x},
\label{eq:weighted-s1}
\end{equation}
其中 $\mathbf{x}_{\sim i},z_i$ 表示仅将第 $i$ 维替换为独立采样 $z_i\sim q_i$。实现中使用 Saltelli 风格的矩阵估计器：由 $q$ 采样基线矩阵 $A$ 与重采样矩阵 $B$，令 $AB_i$ 表示把 $A$ 的第 $i$ 列替换为 $B$ 的第 $i$ 列，则
\begin{equation}
\widehat{S}_i^{q}
=
\frac{\frac{1}{N}\sum_{j=1}^{N}
f(B_j)\left[f(AB_{i,j})-f(A_j)\right]}
{\widehat{D}_q},
\qquad
\widehat{S}_{Ti}^{q}
=
\frac{\frac{1}{2N}\sum_{j=1}^{N}
\left[f(A_j)-f(AB_{i,j})\right]^2}
{\widehat{D}_q}.
\label{eq:weighted-sobol-estimator}
\end{equation}
为保证对照公平，加权版与均匀 Sobol 使用同一个随机森林、相同的重采样规模和随机种子，只改变采样测度。该计算需要拟合代理模型并执行 $N(d+2)$ 次批量预测，因此实现为 StaticPlugin，通过 Redis Queue 异步执行并缓存结果。图~\ref{fig:importance} 展示了插件界面：柱状图按加权一阶效应排序，并可叠加均匀测度下的 Sobol 指数作为对照，从而帮助用户区分``全局均匀重要性''与``当前搜索分布下的实际重要性''。

\begin{figure}[t]
    \centering
    \includegraphics[width=0.86\textwidth]{interface/importance.png}
    \caption{Weighted Sobol 重要性界面。图中同时显示经验采样测度下的加权 Sobol 指数与均匀测度对照，超参数按加权一阶效应排序。}
    \label{fig:importance}
\end{figure}

\subsection{Sampling Density}
\label{sec:sampling-density}

Sampling Density 插件用于检查 Weighted Sobol 识别出的重要维度上的实际搜索行为。对一维连续超参数，插件以 Gaussian KDE 估计成功 trial 的经验密度，并绘制均匀分布基线、样本 rug 和 incumbent 位置；类别超参数则直接统计频率。对二维连续超参数，插件使用分箱直方图或热图近似联合密度，并叠加已评估配置与 incumbent。为避免大规模 run 下的高维 KDE 开销，二维视图默认使用不超过 $80\times 80$ 的网格。插件实现为 DynamicPlugin，只有所选超参数、目标或 budget 变化时才重新计算；若 run 很大，则降低网格分辨率以保证交互响应。图~\ref{fig:density} 展示了 1D/2D 密度视图，可用于判断搜索是充分覆盖、过度集中、贴边受限，还是偏向特定离散取值。

\begin{figure}[t]
    \centering
    \includegraphics[width=0.86\textwidth]{interface/density.png}
    \caption{Sampling Density 界面。1D 视图比较经验边际密度与均匀基线并标记 incumbent；2D 视图展示两个超参数上的联合采样覆盖。}
    \label{fig:density}
\end{figure}

\subsection{Footprint Replay}
\label{sec:footprint-replay}

Configuration Footprint 通常展示终态 MDS 投影，难以回答``优化器何时开始收缩、incumbent 何时迁移''。Footprint Replay 在保留现有 Footprint evaluator 的距离计算和 MDS 投影的基础上，为静态点云补充过程时间轴。具体地，系统先在归一化配置距离矩阵上执行预计算距离的 MDS，并一次性冻结二维坐标，避免回放过程中投影漂移；随后按 trial 的完成时间构建两个序列：一是每个配置的 \texttt{first\_seen} 序列，二是按时间前缀维护的 incumbent 轨迹。MDS、性能表面和上述序列由 StaticPlugin 通过 Redis Queue 一次计算并缓存。

界面中的 trial 滑块被放置在 filter 层而非 input 层。因此，用户拖动滑块时不会重新触发 evaluator，而只是在缓存输出的点云和 incumbent 前缀上做数组切片与重绘。图~\ref{fig:replay} 展示了回放界面：灰色点表示尚未访问或早期访问的区域，彩色点表示当前前缀内已评估配置，折线表示 incumbent 迁移路径。该视图使用户能够区分正常收缩、过早收敛和后期无效扩散，从而为下一步的搜索空间干预提供证据。

\begin{figure}[t]
    \centering
    \includegraphics[width=0.92\textwidth]{interface/replay.png}
    \caption{Footprint Replay 界面。MDS 坐标在计算后冻结，trial 滑块只对缓存结果切片；图中同时显示已访问点、性能颜色和 incumbent 前缀轨迹。}
    \label{fig:replay}
\end{figure}

\subsection{Bottleneck Diagnosis}
\label{sec:Bottleneck Diagnosis}
Bottleneck Diagnosis 在原生 Cost over Time 之上，对给定 budget 与 seed 下成功 trial 的代价曲线做流式低效区间检测与诊断。系统先把 run 转为按完成时间排序的记录序列（聚合全部 seed 时取每配置平均代价，目标统一归一化为越小越好），其上滑动大小 $k=20$ 的窗口（自动收缩至不超过 trial 数的 10\%），计算 incumbent 相对改进率 $\Delta=(B_{\text{before}}-B_{\text{after}})/|B_{\text{before}}|$。若 $\Delta\ge\varepsilon$（$\varepsilon=10^{-3}$）则判定窗口仍在改进；否则对窗口前半与后半的原始代价分布做 Mann–Whitney U 检验，若显著改善（$p<0.05$）归为噪声级进展而非停滞。对真正停滞的窗口计算新颖率 $e$：以最近 200 个新配置最近邻距离的 $q=0.75$ 分位数为自适应阈值 $\tau$，$e$ 为窗口内新配置中归一化最近邻距离不小于 $\tau$ 的比例；$e<\theta_e=0.5$ 判为局部最优停滞，否则判为无效探索。

区段状态机在连续两个同类窗口后才确认停滞区间，并把区间起点锚定在最后一次显著改进处；区间长度达 $3k$ 或超过全部 trial 的 30\% 记为 critical。插件在代价曲线上叠加彩色阴影区间并给出状态横幅：内置按优先级排序的双语建议规则库依据上下文生成可操作建议，如停滞超过全程 50\% 且仍在继续时建议终止重启、局部最优时建议增强探索或重新参数化搜索空间、无效探索时建议加强利用，以及噪声过大时建议提高保真度。图~\ref{fig:bottleneck} 展示了实时识别停滞空间并提出建议的插件输出。

\begin{figure}[t]
    \centering
    \includegraphics[width=0.92\textwidth]{interface/bottleneck.png}
    \caption{Bottleneck Diagnosis 界面。在Cost over time曲线上，基于Mann–Whitney U 检验实时识别并标记低效探索区间，并通过计算参数新颖率e，识别是否是重复型低效探索。}
    \label{fig:bottleneck}
\end{figure}


\subsection{Hypervolume Convergence}
\label{sec:Hypervolume}
Hypervolume Convergence 用 Zitzler 与 Thiele 于1999年提出的超体积（HV）指标度量多目标优化质量。HV 是唯一严格 Pareto 兼容的单 unary 质量指标：前沿不劣化时指标值不会下降，因而无需已知真实前沿即可横向比较 run。给定参考点 $r$ 与非支配解集 $P$，其值为 $P$ 相对 $r$ 所支配的目标空间体积，
$$\mathrm{HV}(P,r)=\Lambda\Big(\bigcup_{p\in P}[p_1,r_1]\times\cdots\times[p_d,r_d]\Big).$$
计算按目标数分级：一维取闭式差；二维将前沿按第一目标升序排序后单次扫描积分 $\mathrm{HV}=\sum_i (x_{i+1}-f_1(p_i))(r_2-f_2(p_i))$，复杂度 $O(n\log n)$；三维用切片法，逐层调用二维精确解；四维及以上采用 WFG 精确算法（While、Bradstreet 与 Barone，2012）按独占体积分解累加，当前沿超过 500 点时改用 $10^5$ 样本的 Monte Carlo 近似以约束运行时间。参考点取各目标联合最差值的 1.1 倍（负值改作除法），全部点再按目标归一化到单位立方，消除量纲差异。

插件沿 trial 提交顺序取至多 400 个检查点增量计算收敛序列，界面绘制多 run 对比曲线（横轴可切换 trial 数、墙钟时间或对数时间），并在汇总表中给出梯形 AUC、首次达到终值 90\% 的时刻 $T_{90}$，以及每个 run 实际命中的算法层级徽标（精确/WFG/Monte Carlo），帮助用户比较收敛速度并识别过早收敛的 run。图~\ref{fig:HyperVolume} 展示了三个run在accuracy–Time 权衡集合上超体积收敛曲线。

\begin{figure}[t]
    \centering
    \includegraphics[width=0.92\textwidth]{interface/HyperVolume.png}
    \caption{HyperVolume Convergence界面。阶梯线表示截至该 trial 的 Pareto 前沿超体积（HV）；HV 越高，表示找到的多目标优化权衡集合整体越好。图中 run\_2 的最终 HV 最高，较早达到自身最终 HV 的 90\%，是这组运行里最终权衡质量和早期进展都较突出的 run。}
    \label{fig:HyperVolume}
\end{figure}


\subsection{Fidelity Replay}
\label{sec:Fidelity Replay}
Budget Correlation 只对相邻预算给出 Spearman 秩相关，无法回答 multi-fidelity 场景下的关键问题：现有预算策略节省了多少计算、又牺牲了多少最优解。Fidelity Replay 据此把既有 trial 历史当作一次 Successive Halving / Hyperband / ASHA 式调度来重放：逐预算层级收集每配置的平均代价（归一化为越小越好），模拟自底向上执行，每层仅将排名前 $1/\eta$ 的活跃配置晋升到下一层；被晋升但下一层缺评估的配置按缺失策略处理——`carry last`（默认，沿用最近已知代价）、`skip`（从阶梯剔除）或 `worst`（按该层最差代价记账）。核心参数 $\eta$ 支持 auto，由预算阶梯自动推断整数几何比，失败时回退为 3。

核心输出是模拟与真实对照的双曲线：实线为各层晋升集合中的最优代价，虚线为该层全部真实评估的最优代价，横轴为对数预算；同一层两线的差距即“模拟调度换算力”所牺牲的 incumbent 质量，而 Active 与 Promoted 的逐层衰减则量化了模拟会在高保真度上省去的评估量（相当于被剪枝的配置数），两者合起来正面回答“省多少、亏多少”。下方逐层 KPI 表列出 Evaluated、Active、Promoted、两线最优值、真实覆盖率与相邻层秩相关，组模式另附成员对照表与 Mean 汇总行；横幅显示 $\eta$（推断/手动）与缺失策略徽标。插件为 DynamicPlugin，参数变化时即时重算。图~\ref{fig:Fidelity} 展示了一组run在各档预算（对应横轴budget）下实际评估配置的最优accuracy（对应纵轴）。

\begin{figure}[t]
    \centering
    \includegraphics[width=0.92\textwidth]{interface/Fidelity.png}
    \caption{Fidelity Replay界面。图中显示，低预算排名与高预算结果已有一定相关性；budget = 80 带来了最后一次可见的最佳 accuracy 提升,可以把 80 作为当前值得评估的候选停止点.}
    \label{fig:Fidelity}
\end{figure}

\section{附录}

\subsection{复现方法}

% =========================================================
\section{References}
% =========================================================

\begin{thebibliography}{99}

\bibitem{bergstra2012random}
J.~Bergstra and Y.~Bengio.
\newblock Random search for hyper-parameter optimization.
\newblock \emph{Journal of Machine Learning Research}, 13:281--305, 2012.

\bibitem{biedenkapp2018cave}
A.~Biedenkapp, H.~F. Bozkurt, T.~Eimer, F.~Hutter, and M.~Lindauer.
\newblock {CAVE}: Configuration assessment, visualization and evaluation.
\newblock In \emph{International Conference on Learning and Intelligent Optimization}, 2018.

\bibitem{bischl2021hpo}
B.~Bischl, M.~Binder, M.~Lang, T.~Pielok, J.~Richter, S.~Coors, A.~Thomas, T.~Ullmann, M.~Becker, A.~Biedenkapp, et~al.
\newblock Hyperparameter optimization: Foundations, algorithms, best practices and open challenges.
\newblock \emph{WIREs Data Mining and Knowledge Discovery}, 13\penalty0 (2):e1440, 2023.

\bibitem{endert2012human}
A.~Endert, W.~Ribarsky, C.~Turkay, B.~L.~W. Wong, I.~Nabney, I.~D. Blanco, and F.~Spaccapietra.
\newblock The state of the art in integrating machine learning into visual analytics.
\newblock \emph{Computer Graphics Forum}, 31\penalty0 (8):2442--2463, 2012.

\bibitem{erickson2020autogluon}
N.~Erickson, J.~Mueller, A.~Shirkov, H.~Zhang, P.~Liu, A.~Zimovets, A.~Smola, and A.~C. Mu.
\newblock {AutoGluon-Tabular}: Robust and accurate {AutoML} for structured data.
\newblock \emph{arXiv preprint arXiv:2003.06505}, 2020.

\bibitem{falkner2018bohb}
S.~Falkner, A.~Klein, and F.~Hutter.
\newblock {BOHB}: Robust and efficient hyperparameter optimization at scale.
\newblock In \emph{International Conference on Machine Learning}, pages 1436--1445, 2018.

\bibitem{frazier2018tutorial}
P.~I. Frazier.
\newblock A tutorial on {Bayesian} optimization.
\newblock \emph{arXiv preprint arXiv:1807.02811}, 2018.

\bibitem{giovanelli2023interactive}
J.~Giovanelli, A.~Tornede, T.~Tornede, and M.~Lindauer.
\newblock Interactive hyperparameter optimization in multi-objective problems via preference learning.
\newblock \emph{arXiv preprint arXiv:2309.03581}, 2023.

\bibitem{holzleitner2021hpo}
M.~Holzleitner, L.~Malyuta, M.~Hutter, and S.~Becker.
\newblock A comparison of hyperparameter optimization methods, discrete acquisition functions, and activation functions for economic forecast ensembles.
\newblock \emph{arXiv preprint arXiv:2111.08563}, 2021.

\bibitem{hutter2014importance}
F.~Hutter, H.~Hoos, and K.~Leyton-Brown.
\newblock An efficient approach for assessing hyperparameter importance.
\newblock In \emph{International Conference on Machine Learning}, pages 754--762, 2014.

\bibitem{hyptertuner2018}
T.~Li, G.~Convertino, W.~Wang, H.~Most, T.~Zajonc, and Y.-H. Tsai.
\newblock {HyperTuner}: Visual analytics for hyperparameter tuning by professionals.
\newblock In \emph{IEEE Workshop on Machine Learning from User Interaction for Visualization and Analytics}, 2018.

\bibitem{jamieson2016nonstochastic}
K.~Jamieson and A.~Talwalkar.
\newblock Non-stochastic best arm identification and hyperparameter optimization.
\newblock In \emph{Artificial Intelligence and Statistics}, pages 486--495, 2016.

\bibitem{jones1998efficient}
D.~R. Jones, M.~Schonlau, and W.~J. Welch.
\newblock Efficient global optimization of expensive black-box functions.
\newblock \emph{Journal of Global Optimization}, 13\penalty0 (4):455--492, 1998.

\bibitem{keim2008visual}
D.~Keim, G.~Andrienko, J.-D. Fekete, C.~G{\"o}rg, J.~Kohlhammer, and G.~Melan{\c c}on.
\newblock Visual analytics: Definition, process, and challenges.
\newblock In \emph{Lecture Notes in Computer Science}, volume 4950, pages 154--175. Springer, 2008.

\bibitem{klein2017fast}
A.~Klein, S.~Falkner, S.~Bartels, P.~Hennig, and F.~Hutter.
\newblock Fast {Bayesian} optimization of machine learning hyperparameters on large datasets.
\newblock In \emph{Artificial Intelligence and Statistics}, pages 528--536, 2017.

\bibitem{kruskal1964mds}
J.~B. Kruskal.
\newblock Multidimensional scaling by optimizing goodness of fit to a nonmetric hypothesis.
\newblock \emph{Psychometrika}, 29\penalty0 (1):1--27, 1964.

\bibitem{kucherenko2012dependent}
S.~Kucherenko, S.~Tarantola, and P.~Annoni.
\newblock Estimation of global sensitivity indices for models with dependent variables.
\newblock \emph{Computer Physics Communications}, 183\penalty0 (4):937--946, 2012.

\bibitem{li2018hyperband}
L.~Li, K.~Jamieson, G.~DeSalvo, A.~Rostamizadeh, and A.~Talwalkar.
\newblock Hyperband: A novel bandit-based approach to hyperparameter optimization.
\newblock \emph{Journal of Machine Learning Research}, 18\penalty0 (185):1--52, 2018.

\bibitem{li2020asha}
L.~Li, K.~Jamieson, A.~Rostamizadeh, E.~Gonina, M.~Hardt, B.~Recht, and A.~Talwalkar.
\newblock A system for massively parallel hyperparameter tuning.
\newblock In \emph{Conference on Machine Learning and Systems}, pages 244--262, 2020.

\bibitem{sach2016progress}
M.~Sach, D.~Radev, T.~Ertl, and S.~Diehl.
\newblock Progress: Supporting visual analysis with task progress.
\newblock In \emph{EuroVis Workshop on Visual Analytics}, 2016.

\bibitem{saltelli2008global}
A.~Saltelli, M.~Ratto, T.~Andres, F.~Campolongo, J.~Cariboni, D.~Gatelli, M.~Saisana, and S.~Tarantola.
\newblock \emph{Global Sensitivity Analysis: The Primer}.
\newblock Wiley, 2008.

\bibitem{saltelli2010variance}
A.~Saltelli, P.~Annoni, I.~Azzini, F.~Campolongo, M.~Ratto, and S.~Tarantola.
\newblock Variance based sensitivity analysis of model output: Design and estimator for the total sensitivity index.
\newblock \emph{Computer Physics Communications}, 181\penalty0 (2):259--270, 2010.

\bibitem{segel2025deepcave}
S.~Segel, H.~Graf, E.~Bergman, K.~Thieme, M.~Wever, A.~Tornede, F.~Hutter, and M.~Lindauer.
\newblock {DeepCAVE}: A visualization and analysis tool for automated machine learning.
\newblock \emph{arXiv preprint arXiv:2512.01810}, 2025.

\bibitem{shahriari2016taking}
B.~Shahriari, K.~Swersky, Z.~Wang, R.~P. Adams, and N.~de~Freitas.
\newblock Taking the human out of the loop: A review of {Bayesian} optimization.
\newblock \emph{Proceedings of the IEEE}, 104\penalty0 (1):148--173, 2016.

\bibitem{snoek2012practical}
J.~Snoek, H.~Larochelle, and R.~P. Adams.
\newblock Practical {Bayesian} optimization of machine learning algorithms.
\newblock In \emph{Advances in Neural Information Processing Systems}, pages 2951--2959, 2012.

\bibitem{srinivas2010gaussian}
N.~Srinivas, A.~Krause, S.~M. Kakade, and M.~W. Seeger.
\newblock {Gaussian} process optimization in the bandit setting: No regret and experimental design.
\newblock In \emph{International Conference on Machine Learning}, pages 1015--1022, 2010.

\bibitem{sobol2001global}
I.~M. Sobol.
\newblock Global sensitivity indices for nonlinear mathematical models and their {Monte Carlo} estimates.
\newblock \emph{Mathematics and Computers in Simulation}, 55\penalty0 (1--3):271--280, 2001.

\bibitem{while2006faster}
L.~While, P.~Hingston, L.~Barone, and T.~Huband.
\newblock A faster algorithm for calculating hypervolume.
\newblock \emph{IEEE Transactions on Evolutionary Computation}, 10\penalty0 (1):29--38, 2006.

\bibitem{zitzler1999spea}
E.~Zitzler and L.~Thiele.
\newblock Multiobjective evolutionary algorithms: A comparative case study and the strength {Pareto} approach.
\newblock \emph{IEEE Transactions on Evolutionary Computation}, 3\penalty0 (4):257--271, 1999.

\end{thebibliography}

\end{document}